% 第四章引言 + 4.1 诱导表示

在第三章的最后两节中，我们只对空间群的表示理论给出了并不完全的叙述。若干论证仅以梗概给出，还有一些论断未加证明。这样处理有几个原因。第一，我们觉得对某些读者来说第三章的叙述已经足够；的确，一个完全严格的论述的细节往往会把读者的注意力从中心议题上引开，使他在晦涩上的所失与在精确上的所得一样多。第二，完全严格的方案将适用于更广的群论背景，不仅对经典空间群有用，对其他应用也有用。因此，用一套不局限于空间群的普遍记号、并使之独立于上一章的那些特殊考虑来陈述完整的群论，似乎是明智的。

因此，本章的目的是以完全普遍的方式处理具有不变子群的群的表示理论；空间群是它的一个特例，其不变子群是三维平移群，因而是 Abel 群。这一补充理论将使我们能够填补上一章的空白，主要是证明为求空间群的表示而概述的方案确实产生全部表示——既不多也不少，也没有与其它表示等价的赝品——同时说明如何由已列成表的小表示得到精确矩阵形式的完全表示。附带地，这一理论还使我们能更容易地讨论固体理论中另外一些有趣的问题，例如表示的实性以及表示之间的内直积。实性之所以重要，是因为它有助于研究晶体中由于时间反演对称而引起的单个粒子或准粒子态的简并（Herring 1937\textit{a}, Wigner 1932）。内直积之所以重要，是因为它有助于确定固体中涉及电子、声子、磁振子等的各种过程的选择定则（Elliott and Loudon 1960）。

\section{诱导表示}

设 $\mathbf{G}$ 是一个群，$\mathbf{K}_{1}$ 是 $\mathbf{G}$ 的一个不一定是 invariant 的子群。例如 $\mathbf{G}$ 可以是一个空间群，而 $\mathbf{K}_{1}$ 是布里渊区中矢量 $\mathbf{k}_{1}$ 的小群。把 $\mathbf{K}_{1}$ 的元素记作 $k_{a}$，并把 $\mathbf{G}$ 写成
\begin{equation*}
\mathbf{G} = \sum_{\alpha}p_{\alpha}\mathbf{K}_{1},
\tag{4.1.1}
\end{equation*}
其中 $p_{\alpha}$ 是左陪集代表，我们假定对一切 $\alpha$ 一次性选定；例如总取 $p_{1} = e$（恒等元），并且每当我们写 $p_{\lambda}$ 时，指的都是式 (4.1.1) 展开中陪集 $p_{\lambda}\mathbf{K}_{1}$ 所对应的那一个陪集代表。

\noindent\textbf{定理 4.1.1}\quad 每个元素 $g \in \mathbf{G}$ 都可以表示成 $g = k_{b}p^{-1}_{\lambda}$，其中 $b$ 与 $\lambda$ 由 $g$ 唯一确定。

因为由刚才从式 (4.1.1) 推得的结果，$g^{-1} = p_{\lambda}k_{s}$，其中 $\lambda$ 与 $s$ 由 $g^{-1}$ 唯一确定。于是 $g = k^{-1}_{s}p^{-1}_{\lambda} = k_{b}p^{-1}_{\lambda}$，其中由于 $\mathbf{K}_{1}$ 是群，$k_{b} = k^{-1}_{s} \in \mathbf{K}_{1}$，定理得证。

现设 $\Omega_{1}$ 是在 $\mathbf{K}_{1}$ 下不可约的、维数为 $d_{j}$ 的矢量空间，取基 $\langle \phi_{r} |$（$r = 1$ 到 $d_{j}$），使得对一切 $k_{a} \in \mathbf{K}_{1}$
\begin{equation*}
k_{a}\phi_{r} = \sum^{d_{j}}_{p=1}\phi_{p}\Gamma^{j}(k_{a})_{pr}.
\tag{4.1.2}
\end{equation*}
把 $k_{a}$ 在表示 $\Gamma^{j}$ 中的特征标记作 $\chi^{j}(k_{a})$。再设 $\Omega$ 表示由 $d_{j}|\mathbf{G}|/|\mathbf{K}_{1}|$ 个函数 $\phi_{\alpha r}$ 张成的矢量空间，其中对一切 $\alpha$ 与 $r$
\begin{equation*}
\phi_{\alpha r} = p_{\alpha}\phi_{r}.
\tag{4.1.3}
\end{equation*}
式 (4.1.3) 是函数恒等式，即对相同的自变量，$\phi_{\alpha r}$ 与 $p_{\alpha}\phi_{r}$ 取相同的值。在第三章 3.7 节定义形如 $\{T_{i} \mid \mathbf{x}_{i}\}\psi_{\mathbf{k}}$ 的函数时，我们用的正是同样的手法。

\noindent\textbf{定理 4.1.2}\quad $\Omega$ 在 $\mathbf{G}$ 下不变。

因为写 $g = p_{\lambda}k_{s}$，则 $g\phi_{\tau r} = p_{\lambda}k_{s}p_{\tau}\phi_{r} = p_{\gamma}k_{t}\phi_{r}$，其中 $\gamma$ 与 $t$ 由 $g$ 与 $\tau$ 唯一确定。于是用式 (4.1.2) 与 (4.1.3) 得
\begin{equation*}
\begin{aligned}
g\phi_{\tau r} &= p_{\gamma}\sum_{p}\phi_{p}\Gamma^{j}(k_{t})_{pr}\\
&= \sum_{p}\phi_{\gamma p}\Gamma^{j}(k_{t})_{pr}.
\end{aligned}
\tag{4.1.4}
\end{equation*}
式 (4.1.4) 定义了所谓 $\Gamma^{j}$ 在 $\mathbf{G}$ 中的\textit{诱导表示}（induced representation），方便起见记作 $\Gamma^{j} \uparrow \mathbf{G}$，其特征标记作 $\chi(g)$。它的维数是 $d_{j}|\mathbf{G}|/|\mathbf{K}_{1}|$。式 (4.1.4) 的含义是：若 $g$ 把基函数 $\phi_{\tau r}$ 变到 $\Omega$ 的另一个基函数上，则表示矩阵的相应元素由 $\mathbf{K}_{1}$ 的表示矩阵给出。

现在应当清楚，空间群的表示事实上就是小群的小表示的诱导表示。这样我们已经完成了任务的一部分：说明如何由列成表的小群小表示得到空间群表示的详细矩阵形式。所涉步骤由式 (4.1.1)--(4.1.4) 概括，而表示的不可约性由下一节给出的判据处理。

\noindent\textbf{定理 4.1.3}\quad 定义 $\Omega_{\alpha}$ 为由 $d_{j}$ 个函数 $\phi_{\alpha r}$（$r = 1$ 到 $d_{j}$）张成的矢量空间。则对任意固定的 $\alpha$，$\Omega_{\alpha}$ 在子群 $\mathbf{K}_{\alpha} = \{p_{\alpha}k_{a}p^{-1}_{\alpha};\ k_{a} \in \mathbf{K}_{1}\}$ 下不可约。

因为利用式 (4.1.2) 与 (4.1.3)，
\begin{equation*}
p_{\alpha}k_{a}p^{-1}_{\alpha}\phi_{\alpha r} = p_{\alpha}k_{a}\phi_{r} = \sum_{s}\phi_{\alpha s}\Gamma^{j}(k_{a})_{sr}.
\tag{4.1.5}
\end{equation*}
式 (4.1.5) 的含义是：$p_{\alpha}k_{a}p^{-1}_{\alpha}$ 在基 $\langle \phi_{\alpha r} |$ 下的代表矩阵是 $\Gamma^{j}(k_{a})$。而矩阵群 $\Gamma^{j}(k_{a})$ 是不可约的（按定义），定理得证。

子群 $\mathbf{K}_{\alpha}$ 可以写成 $p_{\alpha}\mathbf{K}_{1}p^{-1}_{\alpha}$ 的形式。因此每个子群 $\mathbf{K}_{\alpha}$（$\alpha = 1$ 到 $|\mathbf{G}|/|\mathbf{K}_{1}|$）都具有下列性质：

(i) 它与 $\mathbf{K}_{1}$ 同阶。

(ii) 它在 $\mathbf{G}$ 下与 $\mathbf{K}_{1}$ 共轭，从而与集合 $\mathbf{K}_{\alpha}$ 中其他每个成员共轭（因为共轭是等价关系）。

(iii) 只要各子群互不相同，每一个都与分解 (4.1.1) 中的一个且只有一个陪集 $p_{\alpha}\mathbf{K}_{1}$ 一一对应。

我们认出以前遇到过这种情形。因为若把 $\mathbf{G}$ 认同于空间群的同形点群 $\mathbf{F}$，把 $\mathbf{K}_{1}$ 认同于 $\mathbf{k}_{1}$ 的小陪群，则 $\mathbf{K}_{\alpha}$ 可以认同为 $\mathbf{k}_{\alpha}$ 的小陪群，其中 $\mathbf{k}_{\alpha}$ 是 $\mathbf{k}_{1}$ 的星中与陪集 $p_{\alpha}\mathbf{K}_{1}$ 相应的成员。

\noindent\textbf{定理 4.1.4}\quad （Johnston 不可约性判据）\ 把 $g$ 在 $\Omega_{\alpha}$ 中的特征标记作 $\chi_{\alpha}(g)$，于是 $g$ 在 $\Gamma^{j} \uparrow \mathbf{G}$ 中的特征标 $\chi(g) = \sum_{\alpha}\chi_{\alpha}(g)$。那么只要对一切满足 $\alpha \neq \beta$ 的对 $(\alpha, \beta)$ 有
\begin{equation*}
\sum_{t}\chi^{*}_{\alpha}(t)\chi_{\beta}(t) = 0
\tag{4.1.6}
\end{equation*}
（对每一组选定的 $\alpha, \beta$，对 $t$ 的求和取遍 $\mathbf{K}_{\alpha} \cap \mathbf{K}_{\beta}$ 的所有元素），$\Omega$ 就在 $\mathbf{G}$ 下不可约。

Johnston (1960) 给出的这个有用定理的证明如下。

由定理 1.3.8，$\Omega$ 在 $\mathbf{G}$ 下不可约当且仅当
\begin{equation*}
\sum_{g}\chi^{*}(g)\chi(g) = |\mathbf{G}|,
\tag{4.1.7}
\end{equation*}
其中对 $g$ 的求和取遍 $\mathbf{G}$ 的所有元素。也就是说，若
\begin{equation*}
\sum_{\alpha}\sum_{g}|\chi_{\alpha}(g)|^{2} + \sum_{\alpha \neq \beta}\sum_{g}\left\{\chi^{*}_{\alpha}(g)\chi_{\beta}(g)\right\} = |\mathbf{G}|.
\tag{4.1.8}
\end{equation*}
由式 (4.1.4) 显然，若 $g$ 可写成 $p_{\tau}k_{t}p^{-1}_{\tau}$，则
\begin{equation*}
\chi_{\tau}(g) = \chi^{j}(k_{t}),
\tag{4.1.9a}
\end{equation*}
否则
\begin{equation*}
\chi_{\tau}(g) = 0.
\tag{4.1.9b}
\end{equation*}
于是 $\chi_{\tau}(g) = \chi^{j}(k_{t})$ 若 $g \in \mathbf{K}_{\tau}$ 且 $g = p_{\tau}k_{t}p^{-1}_{\tau}$，而对一切不属于 $\mathbf{K}_{\tau}$ 的 $g$ 它为零。由此
\begin{equation*}
\sum_{\alpha}\sum_{g}|\chi_{\alpha}(g)|^{2} = \sum_{\alpha}\sum_{a}|\chi^{j}(k_{a})|^{2} = \sum_{\alpha}|\mathbf{K}_{1}| = |\mathbf{G}|.
\end{equation*}
（由于 $\Gamma^{j}$ 在 $\mathbf{K}_{1}$ 下不可约，定理 1.3.8 成立，并且在对 $a$ 求和时已用到它。）

因此 $\Gamma^{j} \uparrow \mathbf{G}$ 不可约当且仅当
\begin{equation*}
\sum_{\alpha \neq \beta}\sum_{g}\chi^{*}_{\alpha}(g)\chi_{\beta}(g) = 0.
\tag{4.1.10}
\end{equation*}
特别地，若对所有满足 $\alpha \neq \beta$ 的 $\alpha, \beta$ 都有
\begin{equation*}
\sum_{g}\chi^{*}_{\alpha}(g)\chi_{\beta}(g) = 0
\tag{4.1.11}
\end{equation*}
则 $\Gamma^{j} \uparrow \mathbf{G}$ 不可约。

借助式 (4.1.9)，式 (4.1.11) 左端可能出现非零项的唯一情形是 $g$ 同时属于 $\mathbf{K}_{\alpha}$ 与 $\mathbf{K}_{\beta}$。这就给出条件 (4.1.6)，定理得证。

\noindent\textbf{推论}

（i）若 $\mathbf{K}_{1}$ 在 $\mathbf{G}$ 中不变，则对一切 $\alpha$ 有 $\mathbf{K}_{\alpha} = \mathbf{K}_{1}$，条件 (4.1.6) 变为
\begin{equation*}
\sum_{a}\chi^{*}_{\alpha}(k_{a})\chi_{\beta}(k_{a}) = 0
\tag{4.1.12}
\end{equation*}
对一切满足 $\alpha \neq \beta$ 的 $\alpha, \beta$ 成立。若由基 $\langle \phi_{\alpha r} |$（$\alpha = 1$ 到 $|\mathbf{G}|/|\mathbf{K}_{1}|$）张成的 $\mathbf{K}_{1}$ 的不可约表示彼此不等价，则上式满足。

（ii）条件 (4.1.6) 对每一对满足 $\alpha \neq \beta$ 的 $\alpha, \beta$ 涉及对群 $\mathbf{K}_{\alpha} \cap \mathbf{K}_{\beta}$ 全部元素的求和，要求该和为零。我们即将证明：若对 $\mathbf{K}_{\alpha} \cap \mathbf{K}_{\beta}$ 的某个子群 $\mathbf{H}_{\alpha\beta}$ 有
\begin{equation*}
\sum_{h}\chi^{*}_{\alpha}(h)\chi_{\beta}(h) = 0,
\tag{4.1.13}
\end{equation*}
（对 $h$ 的求和取遍 $\mathbf{H}_{\alpha\beta}$ 的所有元素），则该和为零。若能选取 $\mathbf{H}_{\alpha\beta}$ 为对一切对 $\alpha, \beta$ 都相同的 $\mathbf{G}$ 的子群 $\mathbf{H}$，这一判据特别有用。

为了证明式 (4.1.13) 足以蕴含特定对 $\alpha, \beta$ 的式 (4.1.6)，更仔细地考察式 (4.1.13) 的含义：它说的是 $\mathbf{H}_{\alpha\beta}$ 的两个表示——一个以 $\langle \phi_{\alpha r} |$ 为基、另一个以 $\langle \phi_{\beta r} |$ 为基——当限制在 $\mathbf{H}_{\alpha\beta}$ 上时没有共同的分量。

\noindent\textbf{定理 4.1.5}\quad （Frobenius 互反定理）\ 设 $\Gamma$ 是 $\mathbf{G}$ 的表示，特征标为 $\chi^{\Gamma}(g)$；则 $\Gamma$ 在 $\Gamma^{j} \uparrow \mathbf{G}$ 分解为 $\mathbf{G}$ 的不可约表示时出现的次数，等于表示 $\Gamma^{j}$ 在 $\Gamma$ 限制到 $\mathbf{K}_{1}$（记作 $\Gamma \downarrow \mathbf{K}_{1}$）的分解中出现的次数。

定理证明如下。$\Gamma$ 在 $\Gamma^{j} \uparrow \mathbf{G}$ 中出现的次数等于交织数 $(1/|\mathbf{G}|)\sum_{g}\chi(g)\chi^{\Gamma}(g^{-1})$：
\begin{equation*}
= \frac{1}{|\mathbf{G}|}\sum_{g}\sum_{\alpha}\chi_{\alpha}(g)\chi^{\Gamma}(g^{-1})
\tag{4.1.14}
\end{equation*}
\begin{equation*}
= \frac{1}{|\mathbf{G}|}\sum_{a}\sum_{\alpha}\chi^{j}(k_{a})\chi^{\Gamma}(p_{\alpha}k^{-1}_{a}p^{-1}_{\alpha})
\tag{4.1.15}
\end{equation*}
\begin{equation*}
= \frac{1}{|\mathbf{K}_{1}|}\sum_{a}\chi^{j}(k_{a})\chi^{\Gamma}(k^{-1}_{a}),
\tag{4.1.16}
\end{equation*}
这是一个交织数，表示 $\Gamma^{j}$ 在 $\Gamma \downarrow \mathbf{K}_{1}$ 中出现的次数。从 (4.1.14) 到 (4.1.15) 我们用了式 (4.1.9)，从 (4.1.15) 到 (4.1.16) 用了 $k^{-1}_{a}$ 与 $p_{\alpha}k^{-1}_{a}p^{-1}_{\alpha}$ 属于 $\mathbf{K}_{1}$ 的同一类这一事实。
